\usepackage{mathrsfs,bbm}

\numberwithin{equation}{section}

\newtheorem{Theorem}{Theorem}[section]
\newtheorem{Corollary}[Theorem]{Corollary}
\newtheorem{Lemma}[Theorem]{Lemma}
\newtheorem{Proposition}[Theorem]{Proposition}
 { \theoremstyle{definition}
\newtheorem{Definition}[Theorem]{Definition}
\newtheorem{Example}[Theorem]{Example}
\newtheorem{Remark}[Theorem]{Remark}
}

\newcommand{\e}{\mathrm{e}} % exponential e
\newcommand{\N}{\mathbb{N}}
\newcommand{\R}{\mathbb{R}}
\newcommand{\Z}{\mathbb{Z}}
\newcommand{\C}{\mathbb{C}}
\newcommand{\Q}{\mathbb{Q}}
\renewcommand{\P}{\mathbb{P}}
\newcommand{\E}{\mathbb{E}}
\newcommand{\dd}{\mathrm{d}} %integration d
\newcommand{\eps}{\varepsilon}
\newcommand{\la}{\langle}
\newcommand{\ra}{\rangle}

\newcommand{\be}{\begin{equation}}
\newcommand{\ee}{\end{equation}}

\newcommand{\vect}[1]{\boldsymbol{#1}}

\renewcommand{\Re}{\mathrm{Re}\,} %real part
\renewcommand{\Im}{\mathrm{Im}\,} %imaginary part
\DeclareMathOperator{\Int}{Int}
\DeclareMathOperator{\supp}{supp} % support
\DeclareMathOperator{\const}{const}

\newcommand{\set}[1]{{\left \{ #1 \right \}}}

\newcommand{\one}{\mathbbm{1}}
\newcommand{\1}{\mathbbm{1}}

\AddToHook{cmd/maketitle/after}{\markboth{Floreani, Jansen and Wagner}{Continuum unitary--Meixner action}}
